% Chicken & Vege Stir-Fry Bowl --- Stovetop
% Save as: Stovetop/Chicken_Veg_Stir_Fry_Bowl.tex
% Add to main file with: % Chicken & Vege Stir-Fry Bowl --- Stovetop
% Save as: Stovetop/Chicken_Veg_Stir_Fry_Bowl.tex
% Add to main file with: % Chicken & Vege Stir-Fry Bowl --- Stovetop
% Save as: Stovetop/Chicken_Veg_Stir_Fry_Bowl.tex
% Add to main file with: % Chicken & Vege Stir-Fry Bowl --- Stovetop
% Save as: Stovetop/Chicken_Veg_Stir_Fry_Bowl.tex
% Add to main file with: \include{Stovetop/Chicken_Veg_Stir_Fry_Bowl}

\begin{recipebox}{Chicken \& Veg Stir-Fry Bowl}{2}{30 min}

A veg-loaded, lean stir-fry with diced chicken breast, cabbage, bean sprouts
and grated carrot, served with a small portion of brown rice. High protein,
high volume and low calorie density.

\begin{ingredients}
  \item 400g chicken breast, diced
  \item 300g cabbage, shredded
  \item 150g bean sprouts
  \item 1 carrot, grated
  \item 1 carrot, julienned
  \item 200g frozen mixed vegetables
  \item 1 onion, sliced
  \item 2 garlic cloves, minced
  \item 100g brown rice (uncooked)
  \item 1 tbsp cooking oil
  \item \textbf{Sauce:}
  \begin{itemize}[leftmargin=1.5em]
    \item 3 tbsp low-sodium soy sauce
    \item 1 tbsp oyster sauce
    \item 1 tsp sesame oil
    \item 1 tsp cornflour
    \item 2 tbsp water
  \end{itemize}
\end{ingredients}

\begin{method}
  \item Cook the brown rice according to the packet instructions (about 30 min). Set aside --- this is a small side portion, not the base of the bowl.
  \item Shred the cabbage, grate one carrot, julienne the other, slice the onion and mince the garlic. Have the bean sprouts and frozen veg ready --- stir-frying moves fast.
  \item Mix the soy sauce, oyster sauce, sesame oil, cornflour and water in a small bowl. Set aside.
  \item Heat the oil in a wok or large pan over high heat. Add the chicken and cook until browned and cooked through (5--6 min). Remove and set aside.
  \item In the same pan, add the onion and garlic. Stir-fry for 1 min until fragrant.
  \item Add the cabbage, julienned carrot, frozen veg and grated carrot. Stir-fry on high heat for 3--4 min, tossing frequently, until the cabbage softens but still has some bite.
  \item Return the chicken to the pan. Add the bean sprouts and pour over the sauce. Toss for 1--2 min until the sauce thickens and coats everything. Excess liquid should reduce quickly on high heat --- don't let it pool.
  \item Divide into bowls with a small portion of brown rice on the side or underneath.
\end{method}

\begin{notelist}
  \item High veg-to-protein ratio keeps calorie density low --- roughly 350--400 kcal per serve depending on rice portion.
  \item Bean sprouts are added last to keep their crunch.
  \item Swap the frozen mixed veg for broccoli or snow peas if priced well.
  \item Optional add-ins: minced ginger with the garlic, chilli flakes, a beaten egg scrambled through at the end for extra protein, or a squeeze of lime to finish.
\end{notelist}

\begin{variationbox}{Noodle / Pasta Version}
Swap the rice for 150g (dry weight) hokkien or rice noodles, or pasta, tossed through at the end in step 7. Pasta and noodles soak up more liquid than rice, so add a splash of water or an extra tablespoon of soy sauce if it looks dry. Keep the veg quantities the same. Higher in calories than the rice version.
\end{variationbox}

\end{recipebox}


\begin{recipebox}{Chicken \& Veg Stir-Fry Bowl}{2}{30 min}

A veg-loaded, lean stir-fry with diced chicken breast, cabbage, bean sprouts
and grated carrot, served with a small portion of brown rice. High protein,
high volume and low calorie density.

\begin{ingredients}
  \item 400g chicken breast, diced
  \item 300g cabbage, shredded
  \item 150g bean sprouts
  \item 1 carrot, grated
  \item 1 carrot, julienned
  \item 200g frozen mixed vegetables
  \item 1 onion, sliced
  \item 2 garlic cloves, minced
  \item 100g brown rice (uncooked)
  \item 1 tbsp cooking oil
  \item \textbf{Sauce:}
  \begin{itemize}[leftmargin=1.5em]
    \item 3 tbsp low-sodium soy sauce
    \item 1 tbsp oyster sauce
    \item 1 tsp sesame oil
    \item 1 tsp cornflour
    \item 2 tbsp water
  \end{itemize}
\end{ingredients}

\begin{method}
  \item Cook the brown rice according to the packet instructions (about 30 min). Set aside --- this is a small side portion, not the base of the bowl.
  \item Shred the cabbage, grate one carrot, julienne the other, slice the onion and mince the garlic. Have the bean sprouts and frozen veg ready --- stir-frying moves fast.
  \item Mix the soy sauce, oyster sauce, sesame oil, cornflour and water in a small bowl. Set aside.
  \item Heat the oil in a wok or large pan over high heat. Add the chicken and cook until browned and cooked through (5--6 min). Remove and set aside.
  \item In the same pan, add the onion and garlic. Stir-fry for 1 min until fragrant.
  \item Add the cabbage, julienned carrot, frozen veg and grated carrot. Stir-fry on high heat for 3--4 min, tossing frequently, until the cabbage softens but still has some bite.
  \item Return the chicken to the pan. Add the bean sprouts and pour over the sauce. Toss for 1--2 min until the sauce thickens and coats everything. Excess liquid should reduce quickly on high heat --- don't let it pool.
  \item Divide into bowls with a small portion of brown rice on the side or underneath.
\end{method}

\begin{notelist}
  \item High veg-to-protein ratio keeps calorie density low --- roughly 350--400 kcal per serve depending on rice portion.
  \item Bean sprouts are added last to keep their crunch.
  \item Swap the frozen mixed veg for broccoli or snow peas if priced well.
  \item Optional add-ins: minced ginger with the garlic, chilli flakes, a beaten egg scrambled through at the end for extra protein, or a squeeze of lime to finish.
\end{notelist}

\begin{variationbox}{Noodle / Pasta Version}
Swap the rice for 150g (dry weight) hokkien or rice noodles, or pasta, tossed through at the end in step 7. Pasta and noodles soak up more liquid than rice, so add a splash of water or an extra tablespoon of soy sauce if it looks dry. Keep the veg quantities the same. Higher in calories than the rice version.
\end{variationbox}

\end{recipebox}


\begin{recipebox}{Chicken \& Veg Stir-Fry Bowl}{2}{30 min}

A veg-loaded, lean stir-fry with diced chicken breast, cabbage, bean sprouts
and grated carrot, served with a small portion of brown rice. High protein,
high volume and low calorie density.

\begin{ingredients}
  \item 400g chicken breast, diced
  \item 300g cabbage, shredded
  \item 150g bean sprouts
  \item 1 carrot, grated
  \item 1 carrot, julienned
  \item 200g frozen mixed vegetables
  \item 1 onion, sliced
  \item 2 garlic cloves, minced
  \item 100g brown rice (uncooked)
  \item 1 tbsp cooking oil
  \item \textbf{Sauce:}
  \begin{itemize}[leftmargin=1.5em]
    \item 3 tbsp low-sodium soy sauce
    \item 1 tbsp oyster sauce
    \item 1 tsp sesame oil
    \item 1 tsp cornflour
    \item 2 tbsp water
  \end{itemize}
\end{ingredients}

\begin{method}
  \item Cook the brown rice according to the packet instructions (about 30 min). Set aside --- this is a small side portion, not the base of the bowl.
  \item Shred the cabbage, grate one carrot, julienne the other, slice the onion and mince the garlic. Have the bean sprouts and frozen veg ready --- stir-frying moves fast.
  \item Mix the soy sauce, oyster sauce, sesame oil, cornflour and water in a small bowl. Set aside.
  \item Heat the oil in a wok or large pan over high heat. Add the chicken and cook until browned and cooked through (5--6 min). Remove and set aside.
  \item In the same pan, add the onion and garlic. Stir-fry for 1 min until fragrant.
  \item Add the cabbage, julienned carrot, frozen veg and grated carrot. Stir-fry on high heat for 3--4 min, tossing frequently, until the cabbage softens but still has some bite.
  \item Return the chicken to the pan. Add the bean sprouts and pour over the sauce. Toss for 1--2 min until the sauce thickens and coats everything. Excess liquid should reduce quickly on high heat --- don't let it pool.
  \item Divide into bowls with a small portion of brown rice on the side or underneath.
\end{method}

\begin{notelist}
  \item High veg-to-protein ratio keeps calorie density low --- roughly 350--400 kcal per serve depending on rice portion.
  \item Bean sprouts are added last to keep their crunch.
  \item Swap the frozen mixed veg for broccoli or snow peas if priced well.
  \item Optional add-ins: minced ginger with the garlic, chilli flakes, a beaten egg scrambled through at the end for extra protein, or a squeeze of lime to finish.
\end{notelist}

\begin{variationbox}{Noodle / Pasta Version}
Swap the rice for 150g (dry weight) hokkien or rice noodles, or pasta, tossed through at the end in step 7. Pasta and noodles soak up more liquid than rice, so add a splash of water or an extra tablespoon of soy sauce if it looks dry. Keep the veg quantities the same. Higher in calories than the rice version.
\end{variationbox}

\end{recipebox}


\begin{recipebox}{Chicken \& Veg Stir-Fry Bowl}{2}{30 min}

A veg-loaded, lean stir-fry with diced chicken breast, cabbage, bean sprouts
and grated carrot, served with a small portion of brown rice. High protein,
high volume and low calorie density.

\begin{ingredients}
  \item 400g chicken breast, diced
  \item 300g cabbage, shredded
  \item 150g bean sprouts
  \item 1 carrot, grated
  \item 1 carrot, julienned
  \item 200g frozen mixed vegetables
  \item 1 onion, sliced
  \item 2 garlic cloves, minced
  \item 100g brown rice (uncooked)
  \item 1 tbsp cooking oil
  \item \textbf{Sauce:}
  \begin{itemize}[leftmargin=1.5em]
    \item 3 tbsp low-sodium soy sauce
    \item 1 tbsp oyster sauce
    \item 1 tsp sesame oil
    \item 1 tsp cornflour
    \item 2 tbsp water
  \end{itemize}
\end{ingredients}

\begin{method}
  \item Cook the brown rice according to the packet instructions (about 30 min). Set aside --- this is a small side portion, not the base of the bowl.
  \item Shred the cabbage, grate one carrot, julienne the other, slice the onion and mince the garlic. Have the bean sprouts and frozen veg ready --- stir-frying moves fast.
  \item Mix the soy sauce, oyster sauce, sesame oil, cornflour and water in a small bowl. Set aside.
  \item Heat the oil in a wok or large pan over high heat. Add the chicken and cook until browned and cooked through (5--6 min). Remove and set aside.
  \item In the same pan, add the onion and garlic. Stir-fry for 1 min until fragrant.
  \item Add the cabbage, julienned carrot, frozen veg and grated carrot. Stir-fry on high heat for 3--4 min, tossing frequently, until the cabbage softens but still has some bite.
  \item Return the chicken to the pan. Add the bean sprouts and pour over the sauce. Toss for 1--2 min until the sauce thickens and coats everything. Excess liquid should reduce quickly on high heat --- don't let it pool.
  \item Divide into bowls with a small portion of brown rice on the side or underneath.
\end{method}

\begin{notelist}
  \item High veg-to-protein ratio keeps calorie density low --- roughly 350--400 kcal per serve depending on rice portion.
  \item Bean sprouts are added last to keep their crunch.
  \item Swap the frozen mixed veg for broccoli or snow peas if priced well.
  \item Optional add-ins: minced ginger with the garlic, chilli flakes, a beaten egg scrambled through at the end for extra protein, or a squeeze of lime to finish.
\end{notelist}

\begin{variationbox}{Noodle / Pasta Version}
Swap the rice for 150g (dry weight) hokkien or rice noodles, or pasta, tossed through at the end in step 7. Pasta and noodles soak up more liquid than rice, so add a splash of water or an extra tablespoon of soy sauce if it looks dry. Keep the veg quantities the same. Higher in calories than the rice version.
\end{variationbox}

\end{recipebox}
